\documentclass[border=10pt]{standalone}
\usepackage{fontspec}
\setmainfont{Noto Serif}
\newfontfamily\bengalifont{Noto Serif Bengali}[Script=Bengali]
\usepackage[Latin,Bengali]{ucharclasses}
\setTransitionsFor{Bengali}{\begingroup\bengalifont}{\endgroup}
\usepackage{newunicodechar}
\newunicodechar{।}{{\bengalifont\char"0964}}
\newunicodechar{॥}{{\bengalifont\char"0965}}

\usepackage[table]{xcolor}
\usepackage{booktabs,tabularx,array}
\definecolor{bapilrule}{HTML}{1F2937}
\definecolor{bapilstripe}{HTML}{F3F4F6}
\renewcommand{\arraystretch}{1.28}
\setlength{\tabcolsep}{5pt}

\begin{document}
\begin{minipage}{17cm}
{\large\bfseries সারণি ৪৯।}\ {\large নির্দেশক বাজেটের পরিকল্পনা-অনুমান}\par\vspace{7pt}
\small
\rowcolors{2}{bapilstripe}{white}
\begin{tabularx}{\linewidth}{@{}>{\hsize=0.679\hsize\raggedright\arraybackslash}X>{\hsize=1.286\hsize\raggedright\arraybackslash}X>{\hsize=1.036\hsize\raggedright\arraybackslash}X@{}}
\arrayrulecolor{bapilrule}\toprule
\rowcolor{white}\textbf{বিষয়} & \textbf{অনুমান} & \textbf{ভিত্তি} \\
\midrule
বিনিময় হার & প্রতি মার্কিন ডলারে ১২৪ টাকা & ১৯ অক্টোবর ২০২৬-এ বাংলাদেশ ব্যাংকের ভারিত গড় হার ১২৩.৬৫ টাকা [110] \\
মূল জনবল & প্রথম, দ্বিতীয় ও তৃতীয় বছরে ৯, ১৫ ও ১৯ পূর্ণকালীন সমতুল্য & জনবল পরিকল্পনা (সারণি ১৬) \\
জনবলের ব্যয় & প্রতি পূর্ণকালীন সমতুল্যের পূর্ণ মাসিক ব্যয়: পরিচালক ৩.০ লাখ টাকা; ML ও বাংলা NLP প্রকৌশলী ২.০ লাখ; নিরাপত্তা প্রকৌশলী ১.৮ লাখ; নীতি, ডেটা ও প্ল্যাটফর্ম কর্মী ১.৫ লাখ; শাসন ও UX কর্মী ১.২৫ লাখ; সমন্বয়ক ০.৯ লাখ; প্রশিক্ষণ কর্মকর্তা ০.৮ লাখ & দুর্লভ এআই দক্ষতার প্রতিযোগিতায় নবম জাতীয় বেতন স্কেল ২০২৬-এর মূল বেতনের (মাসিক ২০,০০০ থেকে ১,৫৬,০০০ টাকা) ওপরে [109] \\
ভাষা-সম্পাদক ও ফেলো & খণ্ডকালীন সম্পাদক বছরে ১২ থেকে ১৫ লাখ টাকা; মাসিক ৪০,০০০ টাকায় ৬, ১০ ও ১২ জন ফেলো & ফেলোশিপ কর্মসূচি (অনুচ্ছেদ ৬.৬) \\
যন্ত্রপাতি & সাধারণ-উদ্দেশ্য সার্ভার প্রতিটি ১২,০০০ মার্কিন ডলার; GPU ইনফারেন্স সার্ভার (L40S শ্রেণি) প্রতিটি ২৫,০০০ মার্কিন ডলার; সংরক্ষণক্ষমতা বছরে ৬,০০০ থেকে ৮,০০০ মার্কিন ডলার & ২০২৫ সালে একটি L40S কার্ড প্রায় ৭,৫০০ মার্কিন ডলার, সঙ্গে চেসিস, মেমরি ও সংরক্ষণক্ষমতা [111] \\
হোস্টিং ও পরিচালনা & বছরে ২৫ লাখ থেকে ৪৫ লাখ টাকা & ডেটা-সেন্টারের স্থান, বিদ্যুৎ, ব্যান্ডউইথ, লাইসেন্স ও উন্নয়ন-ক্লাউড \\
মডেল অভিযোজনের কম্পিউটিং & যৌথ জাতীয় কম্পিউটিং থেকে বস্তুগত সহায়তা & সারণি ৩১; খসড়া জাতীয় এআই নীতি [88] \\
বেঞ্চমার্ক টীকাকরণ & মাস ৩০-এর মধ্যে ২,০০০ প্রশ্ন, নিষ্পত্তিসহ প্রতিটি ২,৫০০ টাকা & BanglaPolicyQA পরিকল্পনা (অনুচ্ছেদ ৯.৪, সংযুক্তি ট) \\
প্রশিক্ষণ & মাস ৩৬-এর মধ্যে ৫০০ কর্মকর্তা, অংশগ্রহণকারীপ্রতি ১৫,০০০ টাকা, সঙ্গে বিশেষায়িত কোর্স & প্রশিক্ষণ পাঠ্যক্রম (সারণি ১৮) \\
মূল্যায়ন ও আশ্বাস & মধ্যমেয়াদি মূল্যায়ন ২০ লাখ টাকা; চূড়ান্ত মূল্যায়ন ২৫ লাখ টাকা; বার্ষিক অনুপ্রবেশ পরীক্ষা ও রেড-টিম & অধ্যায় ১৩ ও ১৮ \\
পরোক্ষ ব্যয় ও আপৎকালীন সংস্থান & অন্যান্য সব ব্যয়ের ১০\% & আয়োজক প্রতিষ্ঠানের নীতির সঙ্গে মিলিয়ে নিশ্চিত করতে হবে \\
\bottomrule
\end{tabularx}
\par\vspace{5pt}{\footnotesize\itshape ব্যয়-নিরূপণ সমীক্ষার (ডেলিভারেবল D1.3) পরিকল্পনা-অনুমান, দরপত্র নয়। মডেলটি blueprint/\allowbreak{}src/\allowbreak{}budget\_model.py হিসেবে প্রকাশিত, যাতে যেকোনো অনুমান বদলে বাজেট পুনরায় হিসাব করা যায়।}
\par\vspace{3pt}{\scriptsize\color{gray} উৎস: BAPIL কারিগরি ব্লুপ্রিন্ট, সংস্করণ ১.০ (২০২৬)। CC BY 4.0।}
\end{minipage}
\end{document}
